\documentclass[11pt]{article}
\usepackage[margin=1in]{geometry}
\usepackage{amsmath,amssymb,amsthm,hyperref}
\setlength{\emergencystretch}{2em}
\newtheorem{proposition}{Proposition}
\title{The stated unnormalized Mahler bound fails for a cube}
\author{Conjecture 00000003552}
\date{October 8, 2026}
\begin{document}
\maketitle
\noindent\textbf{Result.} With the volume product defined in the source as ordinary volume times polar volume, the symmetric cube in dimension twelve has product
\[
 \frac{4^{12}}{12!}<\frac12\leq\frac{8}{\pi^2}.
\]
This contradicts the stated dimension-independent lower bound. In dimension three, the cube's product is $32/3$, which also differs from the claimed extremal value $8/\pi^2$.

\section*{An admissible body and its actual polar}
Let $n\geq1$ and use standard coordinates and Lebesgue measure on $\mathbb R^n$. Put
\[
 C_n=[-1,1]^n,\qquad
 C_n^\circ=\{y\in\mathbb R^n:\textstyle\sum_{i=1}^n x_i y_i\leq1
                       \text{ for every }x\in C_n\}.
\]
The cube is compact, convex and centrally symmetric, and the origin is an interior point. Thus it is an admissible symmetric convex body. Its polar is exactly
\[
 C_n^\circ=\{y:\textstyle\sum_{i=1}^n|y_i|\leq1\}.
\]
Indeed, for $x\in C_n$ we have $\sum_i x_i y_i\leq\sum_i|y_i|$. Conversely, choose $x_i=1$ when $y_i\geq0$ and $x_i=-1$ otherwise. This vector lies in $C_n$ and gives $\sum_i x_i y_i=\sum_i|y_i|$. Hence the polar is the genuine cross-polytope determined by the defining pairing.

\section*{Lebesgue volumes}
The box product formula gives $|C_n|=2^n$. In each orthant the polar is a coordinate reflection of the standard simplex $\{t_i\geq0,\ \sum_i t_i\leq1\}$, whose volume is $1/n!$. The orthants overlap only on coordinate hyperplanes of measure zero, so
\[
 |C_n^\circ|=\frac{2^n}{n!},\qquad
 |C_n|\,|C_n^\circ|=\frac{4^n}{n!}.
\]
The Lean proof evaluates the same Lebesgue volume directly by Mathlib's volume formula for $\ell^p$ balls, at $p=1$:
\[
 \left|\left\{x:\sum_i |x_i|\leq1\right\}\right|
 =\frac{(2\Gamma(2))^n}{\Gamma(n+1)}=\frac{2^n}{n!}.
\]
No polar volume or volume-product identity is assumed.

\newpage
\section*{The numerical contradiction}
At $n=12$, direct integer arithmetic gives $4^{12}/12!<1/2$. Since $0<\pi\leq4$, we have $\pi^2\leq16$ and therefore $1/2\leq8/\pi^2$.
\begin{proposition}
There is a compact, convex, centrally symmetric body with the origin in its interior whose unnormalized Mahler volume is strictly less than $8/\pi^2$.
\end{proposition}
\begin{proof}
Take $C_{12}$. The verified polar and volume computations give the strict inequality above.
\end{proof}
For the ordinary three-dimensional cube, the same computation gives
\[
 |C_3|\,|C_3^\circ|=\frac{4^3}{3!}=\frac{32}{3}\ne\frac{8}{\pi^2}.
\]
Thus if the source's word ``cube'' is read specifically in dimension three, it still does not attain the displayed claimed value under the source's stated volume convention.

\section*{Exact scope and normalization}
The source explicitly defines Mahler volume as the product of volume and polar volume; no normalization or dimension restriction is supplied. The dimension-twelve example refutes its literal universal lower bound. The dimension-three calculation separately refutes attainment of the displayed value by the spatial cube. It does not assert a lower-bound violation in dimension three.

This does not refute the standard symmetric Mahler conjecture with the dimension-dependent value $4^n/n!$. Nor does it refute a statement about a normalized product. For example, in dimension two the square has unnormalized product $8$, whereas dividing by the squared area $\pi^2$ of the Euclidean unit disk gives $8/\pi^2$. Such a normalization would change the quantity defined in the source. The asymmetric conjunct is not needed for this disproof.

\section*{Formal verification}
The Lean project uses finite real coordinate space with its genuine product Lebesgue measure. It proves cube convexity, compactness, symmetry and origin interior; identifies the polar through the complete universal pairing condition and a sign-vector witness; evaluates the two volumes; and derives both dimension-specific conclusions. The concluding theorem assembles the admissibility conditions and strict lower-bound violation.

Run \texttt{lake build} in \texttt{lean/} using Lean 4.19.0. Mathlib is pinned to
\begin{center}\small\texttt{c44e0c8ee63ca166450922a373c7409c5d26b00b}.\end{center}
The final theorem audits use only standard logical axioms. The proof contains no assumed certificates or proof placeholders.
\end{document}
